%% kapa: exam=1 round=29 year=2018 period=1 form=A subject=civil_law subject_ko=민법 seq=9 number=9
%% source: Q-Net 감정평가사 기출문제 (cst00309) artlSeq=5202808 file=extracted/1008/5202808/2018년 제29회 감정평가사 1차시험 기출문제/2018년 제29회 감정평가사 1차시험 기출문제/1교시 A형.pdf
%% conversion: kordoc 4.24.51 -> tools/gen1.py (auto); review=proofread-sonnet; flags=none
\begin{question}{9}
진의 아닌 의사표시에 관한 설명으로 옳지 않은 것은? (다툼이 있으면 판례에 따름)
\begin{choices}
  \item 사인의 공법행위에는 적용되지 않으므로 공무원의 사직 의사가 외부에 표시된 이상 그 의사는 표시된 대로 효력을 발생한다.
  \item 진의는 특정한 내용의 의사표시를 하려는 생각을 말하는 것이지 표의자가 진정으로 마음에서 바라는 사항을 뜻하는 것은 아니다.
  \item 표의자가 강박에 의하여 어쩔 수 없이 증여의 의사표시를 하였다면 이는 비진의 표시에 해당하지 않는다.
  \item 표의자가 비진의표시임을 이유로 의사표시의 무효를 주장하는 경우, 비진의표시에 해당한다는 사실은 표의자가 증명해야 한다.
  \item 표의자가 비진의표시임을 이유로 의사표시의 무효를 주장하는 경우, 상대방이 자신의 선의ㆍ무과실을 증명해야 한다.
\end{choices}
\end{question}
\answer{5}
